% =============================================================
%  Introduction générale
% =============================================================
\chapter*{Introduction générale}
\addcontentsline{toc}{chapter}{Introduction générale}
\markboth{Introduction générale}{Introduction générale}

Le secteur des équipements réseau et des télécommunications évolue à une
vitesse que peu d'industries ont connue. Les appareils d'aujourd'hui ---
passerelles résidentielles (\textit{home gateways}), décodeurs multimédias
(\textit{set-top boxes}) et routeurs industriels --- tournent sur des
systèmes Linux complets, accessibles à distance via le protocole
\gls{telnet}. Avant chaque mise en production, les firmwares doivent
passer par une phase de validation fonctionnelle rigoureuse~: c'est une
étape qu'on ne peut pas se permettre de bâcler.

C'est dans ce contexte que s'inscrit notre projet de fin d'études, réalisé
au sein de \textbf{Sagemcom Ezzahra}, centre R\&D spécialisé dans la
conception d'équipements \gls{cpe} pour des opérateurs télécoms à
l'international. Le constat de départ était simple~: le processus de tests
Telnet manuel existant prenait trop de temps et laissait trop de place à
l'erreur. L'équipe d'ingénierie avait besoin d'une solution web centralisée,
automatisée, et intégrée dans un pipeline \gls{cicd}.

La réponse que nous proposons s'appelle \textbf{Telnet Test Manager}~: une
application web fullstack construite sur les principes de la \textit{Clean
Architecture}, conteneurisée avec \textbf{Docker}, orchestrée par
\textbf{Kubernetes}, et déployée automatiquement via un pipeline
\textbf{GitHub Actions} auto-hébergé.

\vspace{0.4cm}
\noindent Le rapport s'organise en cinq chapitres~:

\begin{description}[labelwidth=3.5cm,leftmargin=4cm]
  \item[\textbf{Chapitre 1}] contexte général~: présentation de l'organisme
    d'accueil, problématique, objectifs, analyse de l'existant et
    méthodologie Scrum.

  \item[\textbf{Chapitre 2}] analyse des besoins~: acteurs du système,
    besoins fonctionnels et non fonctionnels, diagramme des cas
    d'utilisation.

  \item[\textbf{Chapitre 3}] conception~: architecture Clean Architecture,
    diagramme de classes, modèle de données MongoDB et architecture
    technique globale.

  \item[\textbf{Chapitre 4}] implémentation~: moteur Telnet, API REST,
    WebSocket temps réel, interface React et gestion des rôles \gls{rbac}.

  \item[\textbf{Chapitre 5}] infrastructure DevOps~: conteneurisation
    Docker, orchestration Kubernetes/Helm, pipeline CI/CD GitHub Actions,
    supervision Prometheus/Grafana et tests de validation.
\end{description}

\noindent Le rapport se termine par un bilan des réalisations et quelques
pistes d'évolution pour la solution.
